\section{Corrected Follow-up: Matched Static Attribution}\label{sec:dynamicresults}
The main current comparison is the November follow-up, not historical R2. It contains 60 accepted 500-second runs: six methods in ten trace-matched blocks. Classifiers, threshold 0.5 and held-out pool are unchanged. New policies were trained on training caches and selected by validation before deployment. Historical test outcomes had already been inspected before this follow-up was designed; it is not a new untouched-test experiment.

\subsection{Realized Actions and Feasibility}
In all 5,000 windows per method, Q, DQN and Static-LightLR use LightLR; DeadlineGuard and Static-MedRF use MedRF. Their within-block detection summaries coincide. No method records a deadline-miss window. Thus neither learned policies nor DeadlineGuard demonstrate useful load-conditioned switching here. Historical Tabular-Q TinyDT behavior is not reproduced under the corrected interface. This does not establish global policy optimality.

\subsection{Paired Energy and Detection}
DQN minus Static-LightLR is +0.599 J, paired Student-$t$ 95\% CI $[-4.119,5.317]$ J. Q minus Static-LightLR is +0.246 J $[-4.055,4.547]$ J. Neither contrast resolves energy savings. No equivalence margin was declared, so an interval crossing zero does not establish equivalence.

Cascade minus Static-MedRF is $-12.064$ J $[-16.548,-7.580]$ J. Mean F1 changes by $-0.0004457$ ($-0.0446$ percentage points), recall by $-0.2411$ pp and FPR by $-0.1534$ pp. This is a detection--energy trade-off, not a free improvement. The historical 1.373-pp DeadlineGuard trade-off is not a new-study result.
Tables~\ref{tab:followup} and~\ref{tab:sensitivity} retain method summaries and all registered sensitivity endpoints for the key contrasts.
\begin{table*}[t]\centering\small
\caption{Corrected follow-up: means over ten blocks; energy mean $\pm$ sample SD. Each method has zero deadline-miss windows. No historical experiment is pooled.}\label{tab:followup}
\begin{tabular}{lrrrr}\toprule Method & Energy (J) & F1 & Recall (\%) & FPR (\%)\\\midrule
Cascade & $1480.56\pm9.28$ & 0.99017 & 99.408 & 1.364\\
DQN & $1440.57\pm8.73$ & 0.97633 & 96.680 & 1.351\\
DeadlineGuard & $1492.11\pm10.75$ & 0.99062 & 99.649 & 1.517\\
Static-LightLR & $1439.97\pm9.95$ & 0.97633 & 96.680 & 1.351\\
Static-MedRF & $1492.63\pm10.58$ & 0.99062 & 99.649 & 1.517\\
Q & $1440.21\pm9.57$ & 0.97633 & 96.680 & 1.351\\
\bottomrule\end{tabular}\end{table*}
\begin{table*}[t]\centering\footnotesize
\caption{All registered block sets, nominal paired 95\% intervals. A: all blocks; B: exclude 1; C: exclude 1 and 5. Energy differences are first minus second; adjusted energy subtracts preceding idle power times observed duration.}\label{tab:sensitivity}
\begin{tabular}{llllrr}\toprule Set & Contrast & Endpoint & $n$ & Difference (J) & 95\% CI (J)\\\midrule
A & Cascade--Static-MedRF & Raw & 10 & -12.064 & $[-16.548,-7.580]$\\
B & Cascade--Static-MedRF & Raw & 9 & -11.717 & $[-16.748,-6.686]$\\
C & Cascade--Static-MedRF & Raw & 8 & -11.359 & $[-17.129,-5.588]$\\
A & Cascade--Static-MedRF & Adjusted & 10 & -16.089 & $[-26.719,-5.460]$\\
B & Cascade--Static-MedRF & Adjusted & 9 & -16.412 & $[-28.498,-4.326]$\\
C & Cascade--Static-MedRF & Adjusted & 8 & -14.862 & $[-28.287,-1.438]$\\
A & DQN--Static-LightLR & Raw & 10 & 0.599 & $[-4.119,5.317]$\\
B & DQN--Static-LightLR & Raw & 9 & 1.045 & $[-4.207,6.298]$\\
C & DQN--Static-LightLR & Raw & 8 & 2.376 & $[-2.580,7.333]$\\
A & DQN--Static-LightLR & Adjusted & 10 & 0.974 & $[-9.667,11.615]$\\
B & DQN--Static-LightLR & Adjusted & 9 & 2.117 & $[-9.647,13.882]$\\
C & DQN--Static-LightLR & Adjusted & 8 & 3.775 & $[-9.161,16.711]$\\
A & Q--Static-LightLR & Raw & 10 & 0.246 & $[-4.055,4.547]$\\
B & Q--Static-LightLR & Raw & 9 & 0.489 & $[-4.372,5.350]$\\
C & Q--Static-LightLR & Raw & 8 & 2.235 & $[-0.934,5.404]$\\
A & Q--Static-LightLR & Adjusted & 10 & -4.793 & $[-17.147,7.562]$\\
B & Q--Static-LightLR & Adjusted & 9 & -2.694 & $[-15.693,10.306]$\\
C & Q--Static-LightLR & Adjusted & 8 & 1.593 & $[-8.224,11.410]$\\
\bottomrule\end{tabular}\end{table*}


\subsection{Interruptions and Sensitivity}
Two rejected attempts are preserved: block 1 Static-LightLR exceeded 5\% external CPU, coincident with man-db maintenance; block 5 Cascade's idle failed the 1.5-second monitor-gap limit. Accepted runs were retained; explicitly authorized new IDs recover only the interrupted portion. Both blocks span periods. Idle subtraction does not guarantee elimination of temporal confounding.

The derived monitor waits relative to actual sampling time rather than adding one second after processing, and retains rejected snapshots. CPU accounting, 5\% and 1.5-second limits, service loop and policies are unchanged. The old rejected gap lacks its right snapshot; exact duration is unavailable. Both executed source versions are preserved.

We report A: all ten blocks; B: exclude block 1 ($n=9$); C: exclude blocks 1 and 5 ($n=8$), each with raw and preceding-idle-adjusted results. Registration timestamps are actual campaign-time additions, not backdated original preregistration. All method pairs are retained; nominal intervals are not multiplicity-adjusted and no favourable set is selected.

\subsection{Historical Evidence Boundary}
The original study, P0/P1 and old R1/R2 are not pooled. Old R2 used mismatched arrival timing and post-append queue encoding; its TinyDT behavior is an implementation diagnostic, not an interface-matched algorithm comparison. Constant-coordinate ablations do not establish load/queue/thermal robustness. The historical DQN block-4 excursion lacked the later process-aware gate; no precise external-process joule attribution follows. R3 remains NOT EXECUTED; timing is not controller-only energy. Threshold diagnostics do not replace frozen threshold=0.5 results.
